% Portada de M2.5: una señal sinusoidal y su espectro con tres ventanas.
\documentclass[tikz,border=0pt]{standalone}
\usepackage{resonancia}
\begin{document}
\begin{tikzpicture}
  \clip[rounded corners=6pt] (0,0) rectangle (12,3.6);
  \shade[top color=rAzul!22, bottom color=white] (0,0) rectangle (12,3.6);
  % señal
  \draw[eje] (0.5,1.8) -- (5.2,1.8);
  \draw[medida, grueso, smooth] plot[domain=0.6:5.0, samples=200] (\x, {1.8+0.8*sin(520*(\x-0.6))*(0.5+0.5*sin(41*(\x-0.6)))});
  \node[etiqueta, color=rAzul] at (2.8,3.1) {\small $x(t)$};
  \draw[flecha, color=rGris] (5.5,1.8) -- node[above] {\small DFT} (6.7,1.8);
  % espectros
  \draw[eje] (7.0,0.6) -- (11.7,0.6);
  \draw[color=rNaranja, trazo, smooth] plot[domain=0:1, samples=200] ({7.1+4.4*\x}, {0.65+2.1*exp(-((\x-0.5)/0.035)^2)+0.5*abs(sin(500*(\x-0.5)))/(1+60*abs(\x-0.5))});
  \draw[color=rVerde, trazo, smooth] plot[domain=0:1, samples=200] ({7.1+4.4*\x}, {0.65+1.6*exp(-((\x-0.5)/0.06)^2)});
  \node[etiqueta, color=rNaranja, anchor=west] at (7.1,3.1) {\small sin ventana};
  \node[etiqueta, color=rVerde, anchor=east] at (11.7,3.1) {\small con ventana};
\end{tikzpicture}
\end{document}
